\documentclass[12pt,a4paper]{article}
\usepackage[utf8]{inputenc}
\usepackage[T5]{fontenc}
\usepackage{geometry}
\usepackage{setspace}
\usepackage{graphicx}
\usepackage{booktabs}
\usepackage{float}
\usepackage{amsmath}
\usepackage{tabularx}
\usepackage{longtable}
\usepackage{array}
\geometry{margin=1in}
\onehalfspacing
\setcounter{tocdepth}{3}
\renewcommand{\arraystretch}{1.15}

\title{Báo cáo dự án xử lý ngôn ngữ tự nhiên\\Phân tích tình cảm cho các bài đánh giá thương mại điện tử bằng Kỹ thuật học máy.}
\author{Team SkyDom+}
\date{tháng 3 năm 2026}

\begin{document}

\begin{titlepage}
\centering
\vspace*{1.0cm}
{\LARGE Phân tích tình cảm cho các bài đánh giá thương mại điện tử\par}
\vspace{0.2cm}
{\LARGE sử dụng Kỹ thuật học máy.\par}
\vspace{0.8cm}
{\Large Báo cáo nghiên cứu điển hình\par}
\vspace{0.8cm}
{\large Đội: Đội SkyDom+\par}
\vspace{0.3cm}
{\large Thành viên nhóm:\par}
\vspace{0.4cm}
{\large
\begin{tabular}{@{}>{\raggedright\arraybackslash}p{5.8cm} >{\raggedright\arraybackslash}p{2.7cm}@{}}
Nguyen Nhu Tien Tuan & 22BA13314 \\
Tran Hien Chuong & 22BA13056 \\
Do Gia Bao & 22BA13042 \\
Bui The Anh & 22BA13007 \\
Le Duc An & 22BA13002 \\
Pham Quang Tung & 22BA13320 \\
Khuong Bao Duc & 22BA13077 \\
\end{tabular}\par}
\vfill
{\large Ngày 6 tháng 3 năm 2026\par}
\end{titlepage}

\pagenumbering{roman}

\begin{abstract}
Báo cáo này trình bày một dự án NLP toàn diện để xem xét hiểu biết về dữ liệu Thẻ quà tặng Amazon. Quy trình hiện tại có ba lớp vận hành: phân loại cảm tính với khả năng kiểm soát độ không chắc chắn, trích xuất vấn đề đa nhãn và mở rộng theo ngữ cảnh dựa trên máy biến áp. Tập dữ liệu nguồn có 152.410 bài đánh giá, trong khi chú thích vấn đề chứa 19.000 hàng được hợp nhất và xác thực. Biến thể cảm tính cổ điển đã chọn (V6: từ cộng với ký tự TF-IDF có gắn thẻ phủ định) đạt độ chính xác kiểm tra\_0 là 0,677, thu hồi\_0 là 0,932 và F2\_0 là 0,867. Phần mở rộng máy biến áp đạt độ chính xác kiểm tra\_0 là 0,822, thu hồi\_0 là 0,833 và F2\_0 là 0,831 với phạm vi bao phủ 0,910 ở ngưỡng 0,40/0,60. Trình trích xuất vấn đề đạt kiểm tra micro-F1 là 0,765 cho mô hình từ và 0,763 cho mô hình ký tự từ cộng, với macro-F1 tốt hơn trong thiết lập sau. Để phù hợp chặt chẽ với giáo trình NLP, dự án cũng bao gồm một bộ điểm chuẩn (Hồi quy logistic, các biến thể Naive Bayes, Perceptron, LinearSVM, mất log SGD, logistic ngữ nghĩa SVD, FFNN), khối mô hình ngôn ngữ N-gram, đường cơ sở RNN/LSTM, tạo phẩm thăm dò MLM và đường cơ sở kiểu nhắc nhở LLM. Ma trận phù hợp với khóa học được cập nhật được tính toán từ các tạo phẩm được tạo ra báo cáo mức độ bao phủ giáo trình ước tính là 94,3\%.
 B?n c?p nh?t Level 3 trong m? ngu?n hi?n t?i b? sung t?m class-weight theo t?ng label, calibration x?c su?t cho m? h?nh issue v? command ablation classic; ? l?n ch?y m?i nh?t, ablation ch?n V7 l? bi?n th? c? trade-off negative-first t?t nh?t.
\end{abstract}

\tableofcontents
\newpage
\pagenumbering{arabic}

\section{Giới thiệu}
\subsection{Bối cảnh}
Ngôn ngữ tự nhiên trong đánh giá của người dùng khá ồn ào, ngắn gọn và hỗn tạp. Một đánh giá có thể bao gồm khen ngợi, phàn nàn, tiếng lóng, lỗi đánh máy và các từ miền trong cùng một câu. Điều này làm cho việc đánh giá hiểu được một vấn đề NLP thực tế trong đó cả phương pháp cổ điển và hiện đại đều hữu ích \cite{jurafsky2026,turney2002}.

\subsection{Tuyên bố vấn đề}
Nghiên cứu này tập trung vào các đánh giá về Thẻ quà tặng Amazon. Mục tiêu không chỉ giới hạn ở việc phân biệt cảm xúc tích cực và tiêu cực; đó là xây dựng một quy trình xem xét-hiểu biết để hỗ trợ các quyết định nhận thức được sự không chắc chắn và định tuyến ở cấp độ vấn đề cho các hoạt động dịch vụ hạ nguồn.

\subsection{Mục tiêu và phạm vi}
Mục tiêu của dự án là:
\begin{enumerate}
\item Xây dựng bộ phân loại cảm tính có khả năng thu hồi cao, nhận thức rủi ro để phân loại phản hồi tiêu cực.
\item Trích xuất các danh mục vấn đề cụ thể bằng cách sử dụng phân loại văn bản nhiều nhãn.
\item So sánh các mô hình thưa thớt cổ điển với mô hình theo ngữ cảnh dựa trên máy biến áp.
\item Mở rộng phạm vi bảo hiểm nhân tạo để phù hợp với các chủ đề khóa học NLP cốt lõi và nâng cao.
\item Giữ tất cả các kết quả đầu ra có thể tái tạo từ các lệnh và tệp số liệu đã lưu.
\end{enumerate}

\subsection{Những gì chúng tôi đã làm trong dự án này}
Việc triển khai được cấu trúc như một đường dẫn NLP từ đầu đến cuối. Nó bao gồm chuẩn hóa văn bản, mô hình hóa tính năng từ vựng và từ phụ, xử lý phủ định và tương phản, đào tạo có giám sát, kiểm soát độ không chắc chắn, trích xuất vấn đề, mở rộng biến áp, tổng hợp điểm chuẩn và tích hợp trình diễn.

\subsection{Đóng góp chính}
Những đóng góp chính của trạng thái kho lưu trữ hiện tại là:
\begin{enumerate}
\item Một quy trình cảm tính có thể tái tạo với khả năng thu hồi cấp 0 cao và phân tích đánh đổi ngưỡng rõ ràng.
\item Mô-đun trích xuất vấn đề đa nhãn với khả năng điều chỉnh ngưỡng trên mỗi nhãn và số liệu trên mỗi nhãn.
\item Một phần mở rộng máy biến áp với các số liệu dự đoán có chọn lọc và đánh giá trường hợp khó.
\item Gói nâng cấp giáo trình bao gồm NB, Perceptron, FFNN, N-gram LM, RNN/LSTM, đầu dò MLM và đường cơ sở nhắc nhở LLM.
\item Một bộ tài liệu bao gồm ma trận đóng góp, mẫu cổng chất lượng, bảng điểm, bài kiểm tra và tập lệnh trình diễn.
\end{enumerate}

\subsection{Triển khai thực tế}
Dự án cung cấp cả đường dẫn thực thi CLI và giao diện người dùng nhẹ. Các lệnh cốt lõi được ghi lại trong \texttt{README.md}, kết quả đầu ra được lưu trữ trong \texttt{results/} và kiểm tra khói xác thực các đường dẫn thực thi quan trọng. Cấu hình này hỗ trợ khả năng tái tạo trong cả trình bày học thuật và cài đặt thực hiện lại thực tế.

\subsection{Cấu trúc báo cáo}
Báo cáo tuân theo cấu trúc cấp cao tương tự được sử dụng trong \texttt{GP.pdf}: khung vấn đề, phương pháp luận, thiết lập thử nghiệm, kết quả định lượng, bằng chứng quy trình làm việc, kế hoạch demo, những hạn chế và kết luận.

\section{Bối cảnh và các tác phẩm liên quan}
\subsection{Đường cơ sở phân loại cảm xúc và văn bản}
Các mô hình cảm tính cổ điển sử dụng các tính năng n-gram thưa thớt và bộ phân loại tuyến tính vẫn là các đường cơ sở được thiết lập tốt cho văn bản ồn ào do người dùng tạo \cite{wang2012,manning2008ir}. Chúng có hiệu quả về mặt tính toán, có thể giải thích được thông qua trọng số đặc trưng và phù hợp để kiểm soát rủi ro dựa trên ngưỡng.

\subsection{Phủ định và tương phản trong tình cảm}
Phạm vi phủ định là một thách thức đã biết trong phân tích cảm tính \cite{wiegand2010}. Trong văn bản đánh giá, các mẫu như \textit{không tệ} và các dạng tương phản như \textit{tốt nhưng giao hàng muộn} có thể lật ngược ranh giới quyết định. Điều này thúc đẩy việc gắn thẻ phủ định rõ ràng và thiết kế tính năng nhận biết mệnh đề.

\subsection{Học nhiều nhãn để gắn thẻ vấn đề}
Trích xuất vấn đề là một vấn đề có nhiều nhãn vì một bài đánh giá có thể đề cập đến nhiều danh mục khiếu nại \cite{tsoumakas2011,zhang2014}. Thiết lập một so với phần còn lại với hiệu chuẩn ngưỡng trên mỗi nhãn rất thiết thực đối với việc phân phối nhãn không cân bằng.

\subsection{Mô hình hóa theo ngữ cảnh và phần mở rộng biến áp}
Bộ mã hóa dựa trên máy biến áp nắm bắt ngữ nghĩa theo ngữ cảnh và có thể cải thiện độ ổn định của văn bản ngắn thành ngữ \cite{vaswani2017,devlin2019}. Tuy nhiên, trong môi trường nhận thức được rủi ro, các ngưỡng chọn lọc và sự đánh đổi mục tiêu cụ thể của doanh nghiệp vẫn rất cần thiết.

\section{Định nghĩa vấn đề NLP}
\subsection{Nhiệm vụ A: Hiểu biết về tình cảm}
Nhiệm vụ A là phân loại cảm xúc nhị phân từ văn bản đánh giá. Xếp hạng 1 và 2 được ánh xạ tới lớp phủ định. Xếp hạng 4 và 5 được ánh xạ tới lớp tích cực. Xếp hạng có giá trị 3 được sử dụng để phân tích độ không chắc chắn và không nằm trong chương trình đào tạo được giám sát chặt chẽ \cite{turney2002,wang2012}.

Chính sách quyết định không đặt tính chính xác lên hàng đầu. Trong quá trình phân loại, kết quả âm tính giả trên các đánh giá tiêu cực rất tốn kém. Do đó, việc lựa chọn mô hình tuân theo phương pháp phỏng đoán phủ định đầu tiên: ưu tiên thu hồi cho lớp 0, sau đó là F2 cho lớp 0, sau đó là độ chính xác cho lớp 0.

\subsection{Nhiệm vụ B: Trích xuất vấn đề}
Nhiệm vụ B là trích xuất vấn đề đa nhãn. Một bài đánh giá có thể chứa nhiều vấn đề nên việc gắn thẻ một nhãn là không đủ \cite{tsoumakas2011,zhang2014}. Bộ nhãn đầu ra trong dự án này có chín loại vấn đề và cho phép gán nhiều nhãn trên mỗi hàng.

Nhiệm vụ này ánh xạ văn bản đánh giá phi cấu trúc tới các tín hiệu vấn đề có cấu trúc. Nó cho phép chuyển đổi từ kết quả đầu ra tình cảm chung sang các danh mục hành động vận hành.

\subsection{Nhiệm vụ C: Mở rộng bối cảnh sâu}
Nhiệm vụ C là tinh chỉnh máy biến áp dưới dạng phần mở rộng NLP. Phần mở rộng này nhằm mục đích kiểm tra mô hình hóa ngôn ngữ theo ngữ cảnh sâu sắc đối với các cụm từ ngắn và mẫu phủ định khó. Nó cũng cung cấp sự so sánh trực tiếp với các biểu diễn thưa thớt cổ điển \cite{vaswani2017,devlin2019}.

\section{Đặc điểm dữ liệu và ngôn ngữ}
\subsection{Nguồn dữ liệu và quy mô}
Tệp nguồn chứa 152.410 đánh giá với các trường văn bản và xếp hạng. Bảng \ref{tab:nlp_rating_distribution} hiển thị phân bổ xếp hạng đầy đủ.

\begin{table}[H]
\centering
\caption{Phân phối xếp hạng trong dữ liệu nguồn}
\label{tab:nlp_rating_distribution}
\begin{tabular}{lr}
\toprule
Rating & Number of reviews \\
\midrule
1 & 12,326 \\
2 & 1,873 \\
3 & 3,271 \\
4 & 6,692 \\
5 & 128,248 \\
\bottomrule
\end{tabular}
\end{table}

Nhãn tình cảm mạnh mẽ bao gồm xếp hạng 1, 2, 4 và 5. Điều này mang lại 149.139 hàng. Lớp phủ định có 14.199 hàng và lớp tích cực có 134.940 hàng. Tỷ lệ lớp là khoảng 9,5 đến 1. Sự mất cân bằng này rất quan trọng đối với các lựa chọn lập mô hình NLP và đánh giá tập trung vào thu hồi chính xác \cite{saito2015}.

\subsection{Tiếng ồn ngôn ngữ được quan sát}
Ngôn ngữ đánh giá chứa các dạng nhiễu NLP phổ biến. Tập dữ liệu bao gồm các dạng ngắn như gr8, thx, idk và các cụm từ không chính thức. Nó bao gồm nhiều thái cực khác nhau trong một câu, chẳng hạn như giao hàng tốt nhưng trễ. Nó bao gồm các cụm từ phủ định rõ ràng như không tệ và không tốt.

Tập dữ liệu cũng có các từ dành riêng cho tên miền dành cho thẻ quà tặng, đổi mã, kích hoạt, giao hàng, hỗ trợ và hoàn tiền. Những mô hình này thúc đẩy cả việc lập bản đồ từ vựng và học tập dựa trên dữ liệu.

\subsection{Dữ liệu nhãn vấn đề}
Để trích xuất vấn đề, tập dữ liệu được gắn nhãn thủ công có 19.000 hàng. Bảng \ref{tab:nlp_issue_frequency} hiển thị tần suất nhãn.

\begin{table}[H]
\centering
\caption{Tần suất nhãn phát hành trong bộ chú thích NLP thủ công}
\label{tab:nlp_issue_frequency}
\begin{tabular}{lrr}
\toprule
Label & Count & Share \\
\midrule
delivery\_shipping & 972 & 0.0512 \\
redemption\_activation & 284 & 0.0149 \\
product\_quality & 1,251 & 0.0658 \\
customer\_service & 182 & 0.0096 \\
refund\_return & 93 & 0.0049 \\
usability & 873 & 0.0459 \\
value\_price & 558 & 0.0294 \\
fraud\_scam & 144 & 0.0076 \\
other & 15,960 & 0.8400 \\
\bottomrule
\end{tabular}
\end{table}

Nhãn trung bình cho mỗi bài đánh giá là 1,069. Điều này hỗ trợ lập mô hình nhiều nhãn vì nhiều hàng không phải là văn bản một khái niệm.

\section{Thiết kế xử lý văn bản}
\subsection{Chuẩn hóa}
Quy trình chuẩn hóa cốt lõi sử dụng chuyển đổi chữ thường, xóa URL, dọn dẹp dấu câu và chuẩn hóa khoảng trắng. Điều này làm giảm sự thưa thớt của mã thông báo trong khi vẫn giữ được các từ cảm xúc cốt lõi.

Chuẩn hóa chữ viết tắt tùy chọn ánh xạ các dạng ngắn phổ biến thành dạng đầy đủ. Điều này bao gồm các ánh xạ như gr8 thành tuyệt vời, thx thành cảm ơn, u tới bạn và idk tới tôi không biết.

\subsection{Xử lý phủ định}
Phủ định là một thách thức quan trọng của NLP trong các bài đánh giá ngắn. Dự án sử dụng phương pháp gắn thẻ phủ định. Sau tín hiệu phủ định, một số mã thông báo tiếp theo sẽ nhận được tiền tố NOT\_. Kích thước cửa sổ theo mặc định là ba mã thông báo \cite{wiegand2010}.

Phương pháp này nhẹ về mặt tính toán và hiệu quả đối với các mô hình thưa thớt. Nó giúp tách biệt tốt khỏi NOT\_good và xấu khỏi NOT\_bad trong không gian tính năng.

\subsection{Xử lý độ tương phản}
Đường dẫn cũng mô hình hóa độ tương phản diễn ngôn đơn giản. Khi các dấu hiệu như but, tuy nhiên, mặc dù, tuy nhiên, và chưa xuất hiện, văn bản có thể được chia thành các mệnh đề trái và phải. Sau đó, các tính năng mệnh đề riêng biệt và cờ đánh dấu sẽ được thêm vào.

Thiết kế này nắm bắt một mẫu đánh giá chung trong đó mệnh đề thứ hai mang ý kiến cuối cùng của người dùng.

\subsection{Mô hình cấp độ ký tự}
Các đặc điểm ký tự từ 3 đến 5 gam được sử dụng làm sự kết hợp tùy chọn với các đặc điểm của từ. Điều này giúp khắc phục lỗi đánh máy, tiếng lóng ngắn và các biến thể chính tả nhỏ. Trong các tác vụ NLP trên văn bản của người dùng, điều này thường cải thiện độ mạnh mẽ \cite{joulin2017}.

\section{Kỹ thuật tính năng cho NLP}
\subsection{Biểu diễn thưa thớt cổ điển}
Biểu diễn chính là TF-IDF với từ n-gram. Đường cơ sở sử dụng 1 đến 2 gam. Các biến thể khác kiểm tra từ 1 đến 3 gam, thẻ phủ định, mệnh đề tương phản, gam ký tự và đặc điểm từ vựng \cite{manning2008ir,wang2012}.

Bảng \ref{tab:nlp_variant_design} tóm tắt các biến thể ngữ cảnh được thử nghiệm.

\begin{table}[H]
\centering
\caption{Các biến thể của tính năng ngữ cảnh NLP}
\label{tab:nlp_variant_design}
\begin{tabular}{>{\raggedright\arraybackslash}p{1.2cm} >{\raggedright\arraybackslash}p{12cm}}
\toprule
Biến thể & Thiết kế đặc trưng \\
\midrule
V0 & Baseline word TF-IDF with n-gram range 1 to 2 \\
V1 & Word TF-IDF with n-gram range 1 to 3 \\
V2 & Word 1 to 2 plus negation tagging \\
V3 & Word 1 to 2 plus clause split with contrast marker \\
V4 & Word 1 to 2 plus negation tagging plus clause split \\
V5 & Word 1 to 2 plus character 3 to 5 grams \\
V6 & Word 1 to 2 plus character 3 to 5 grams plus negation tagging \\
V7 & Word 1 to 2 plus negation tagging plus sentiment lexicon block \\
\bottomrule
\end{tabular}
\end{table}

\subsection{Khối đặc trưng từ vựng}
Dự án bao gồm khối từ vựng cảm tính dựa trên AFINN trong một biến thể. Nó tính điểm thô, điểm nhận biết phủ định, điểm nhận biết cường độ và điểm có trọng số tương phản. Điều này bổ sung thêm từ vựng rõ ràng trước biểu diễn thưa thớt \cite{nielsen2011}.

Phần này hữu ích cho việc học khóa học NLP vì nó kết nối các phương pháp từ vựng cổ điển với việc học có giám sát hiện đại trong một quy trình.

\section{Phương pháp mô hình hóa}
\subsection{Bộ phân loại tình cảm}
Công cụ phân loại tình cảm chính là Hồi quy logistic. Lựa chọn tính năng Chi bình phương được áp dụng với các giá trị K cao nhất từ 1.000 đến 10.000. Cài đặt trọng lượng lớp không bao gồm các biến thể không có, cân bằng và có trọng số \cite{forman2003}.

Các mô hình cảm tính bổ sung bao gồm Hồi quy logistic L1 cho hành vi nhúng thưa thớt, Cây quyết định và Rừng ngẫu nhiên.

\subsection{Mô hình trích xuất vấn đề}
Trích xuất vấn đề sử dụng đào tạo một so với phần còn lại trên chín nhãn. Mô hình mặc định là Hồi quy logistic với trọng số lớp cân bằng. Đường cơ sở LinearSVM cũng có sẵn \cite{rifkin2004}.

Việc điều chỉnh ngưỡng được thực hiện trên mỗi nhãn trên dữ liệu xác thực. Điều này rất quan trọng trong NLP nhiều nhãn vì một ngưỡng chung thường không đủ cho tất cả các nhãn.

\subsection{Máy biến áp mở rộng}
Phần mở rộng máy biến áp tinh chỉnh bộ phân loại trình tự với tổn thất theo trọng số lớp. Nó sử dụng văn bản rõ ràng và đánh giá với cùng một khung ngưỡng và độ không chắc chắn để so sánh trực tiếp \cite{devlin2019}.

Tiện ích mở rộng này cung cấp ngữ nghĩa theo ngữ cảnh sâu hơn và giúp phân tích các cụm từ ngắn khó hiểu trong đó các mô hình thưa thớt có thể thất bại.

\section{Thiết lập thử nghiệm}
\subsection{Phân chia và giao thức}
Tất cả các thử nghiệm đều sử dụng hạt giống 42. Các nhãn tình cảm mạnh mẽ được phân chia theo tỷ lệ đào tạo, xác thực và kiểm tra 70, 10 và 20. Đánh giá ba sao được giữ để phân tích độ không chắc chắn.

Việc trích xuất vấn đề sử dụng 19.000 hàng được gắn nhãn với kích thước phân chia là 13.300 cho tập huấn, 1.900 cho xác thực và 3.800 cho kiểm tra trong các lần chạy đã lưu.

\subsection{Số liệu chính}
Đối với cảm tính, các chỉ số chính là độ chính xác\_0, mức thu hồi\_0, F1\_0 và F2\_0, cùng với độ chính xác đầy đủ và F1. Do bộ nhãn mạnh bị mất cân bằng nên sự đánh đổi mô hình được diễn giải chủ yếu thông qua hành vi thu hồi chính xác thay vì các chế độ xem chỉ ROC \cite{saito2015}. Các số liệu chọn lọc trong điều kiện không chắc chắn bao gồm phạm vi bao phủ, tỷ lệ không chắc chắn, độ chính xác có chọn lọc\_0, thu hồi có chọn lọc\_0, F2\_0 chọn lọc và tỷ lệ âm tính giả trên các mẫu được bao phủ.

Để trích xuất vấn đề, các số liệu chính là micro-F1, macro-F1, độ chính xác của tập hợp con, tổn thất hamming và độ chính xác trên mỗi nhãn, thu hồi và F1.

\subsection{Quy tắc bất định}
Chính sách không chắc chắn sử dụng ngưỡng thấp và ngưỡng cao đối với xác suất dương. Nó tuân theo thiết lập bỏ qua phân loại có chọn lọc và có thể trả về mức độ không chắc chắn ở phạm vi trung bình \cite{geifman2017}. Chúng tôi áp dụng phương pháp phỏng đoán bổ sung trong quy trình cổ điển để trả về kết quả không chắc chắn đối với các trường hợp văn bản quá ngắn hoặc thưa thớt.

Chính sách này cải thiện độ tin cậy triển khai vì văn bản có ít thông tin không bị buộc phải đưa ra quyết định khó khăn.

\section{Quy trình làm việc nhóm và đảm bảo chất lượng}
\subsection{Phân rã mô-đun và trách nhiệm}
Để duy trì khả năng mở rộng cho một nhóm gồm bảy thành viên, việc triển khai được phân chia thành các mô-đun kỹ thuật với các yêu cầu về quyền sở hữu và kiểm tra chéo rõ ràng. Bảng \ref{tab:nlp_team_modules} tóm tắt cấu trúc được sử dụng trong quá trình phát triển và chuẩn bị trình bày cuối cùng.

\begin{table}[H]
\centering
\caption{Phân chia mô-đun nhóm và chịu trách nhiệm về chất lượng}
\label{tab:nlp_team_modules}
\begin{tabular}{>{\raggedright\arraybackslash}p{3.2cm} >{\raggedright\arraybackslash}p{4.8cm} >{\raggedright\arraybackslash}p{3.8cm} >{\raggedright\arraybackslash}p{2.6cm}}
\toprule
Module & Main responsibility & Cross-check by non-owner & Main artifact path \\
\midrule
Sentiment core & cleaning, TF-IDF features, Chi-square, logistic and uncertainty policy & rerun from raw JSONL and verify metric consistency & \texttt{results/dm2\_steps/} \\
Issue extraction & annotation schema validation, natural ID template generation, OVR training, per-label thresholds & verify deterministic template order and per-label F1 deltas after regeneration & \texttt{results/issue\_steps/} \\
Transformer extension & fine-tuning, threshold sweep, hard-case comparison & compare against classic under the same evaluation protocol & \texttt{results/nlp\_ext/} \\
Reporting and demonstration & report writing, figure integration, demonstration rehearsal & final end-to-end rerun and Q\&A preparation & \texttt{results/reports/} \\
\bottomrule
\end{tabular}
\end{table}

\subsection{Hiểu biết chéo, đánh giá và tái tạo}
Quy trình công việc được thiết kế sao cho mỗi mô-đun quan trọng có thể được thực thi bởi ít nhất một thành viên không phải là tác giả gốc. Đối với mỗi lần chuyển giao mô-đun, thành viên nhận sẽ chạy lại quy trình từ các lệnh đã lưu và kiểm tra ba cổng: không có rò rỉ kiểm tra huấn luyện, hành vi hạt giống ngẫu nhiên cố định và tính nhất quán giữa các giá trị bảng và tệp CSV đã xuất.

Để kiểm soát chất lượng mã, chúng tôi sử dụng chính sách phân nhánh và đánh giá với một đánh giá không phải của tác giả cho mỗi thay đổi đáng kể. Danh sách kiểm tra đánh giá tập trung vào độ chính xác của số liệu, tính toàn vẹn phân tách và khả năng tái tạo. Quá trình này phù hợp với các yêu cầu của khóa học về phân chia mô-đun, chia sẻ hiểu biết và thực hành đánh giá chéo. Bản cập nhật mới nhất ở mô-đun issue cũng bổ sung cơ chế sắp xếp ID tự nhiên cho template gán nhãn và khóa hành vi này bằng test hồi quy chuyên biệt.

\subsection{Gói tài liệu bằng chứng}
Để làm cho bằng chứng làm việc nhóm có thể được kiểm tra trong bản trình bày cuối cùng, gói bằng chứng bao gồm: ma trận vai trò/mô-đun, nhật ký lệnh để chạy lại, tạo phẩm CSV đầu ra của mô hình và ghi chú đánh giá được liên kết với các mô-đun đã thay đổi. Tài liệu này bổ sung cho các kết quả mô hình định lượng và chứng minh chất lượng quy trình bên cạnh điểm số cuối cùng.

\section{Gói nâng cấp giáo trình không có dữ liệu mới}
Dự án hiện bao gồm một gói bổ sung giúp cải thiện sự liên kết với các chủ đề khóa học NLP mà không cần thêm bất kỳ dữ liệu được gắn nhãn mới nào. Gói này sử dụng lại cùng một chuỗi phân tách, xác thực và kiểm tra, đồng thời giữ nguyên chính sách dọn dẹp văn bản.

Khối nâng cấp đầu tiên là bộ điểm chuẩn cổ điển toàn diện. Nó bao gồm Hồi quy logistic, Bayes Naive đa thức, Bayes Naive bổ sung, Perceptron, LinearSVM, SGD với tính năng mất log, Hồi quy logistic ngữ nghĩa trên vectơ SVD và MLP chuyển tiếp nguồn cấp dữ liệu trên cùng vectơ SVD. Kích thước đào tạo được giới hạn ở 60.000 hàng trong khối này để kiểm soát thời gian chạy trong khi vẫn duy trì tỷ lệ lớp.

Khối nâng cấp thứ hai là khối mô hình ngôn ngữ. Nó đào tạo các mô hình ngôn ngữ unigram và bigram với khả năng làm mịn add-k và báo cáo sự phức tạp trong quá trình phân tách xác thực và kiểm tra. Điều này cung cấp bằng chứng trực tiếp cho chủ đề mô hình ngôn ngữ N-gram trong khóa học \cite{jurafsky2026,manning2008ir}.

Khối nâng cấp thứ ba bổ sung thêm các tạo phẩm theo trình tự và mô hình ngôn ngữ hiện đại, bao gồm đường cơ sở cảm tính RNN/LSTM, báo cáo thăm dò MLM và đường cơ sở ngữ nghĩa kiểu lời nhắc LLM trên cùng một phân chia cảm tính.

Khối cuối cùng xây dựng ma trận bao phủ chủ đề từ các tạo phẩm của dự án. Ma trận sử dụng điểm số chủ đề rõ ràng và báo cáo điểm tổng thể về nội dung giáo trình để thảo luận bào chữa.

\section{Kết quả cho Mô hình tình cảm NLP cổ điển}
\subsection{Lựa chọn đường cơ sở bởi K}
Quét đường cơ sở trên V0 hiển thị cài đặt tốt nhất cho mỗi K theo quy tắc phủ định đầu tiên. Bảng \ref{tab:nlp_chi2_best_per_k} báo cáo các hàng đã chọn từ quá trình xác thực.

\begin{table}[H]
\centering
\caption{Kết quả tình cảm cơ bản tốt nhất trên mỗi K khi xác thực}
\label{tab:nlp_chi2_best_per_k}
\begin{tabular}{lcccccc}
\toprule
K & Trọng số lớp & Precision\_0 & Recall\_0 & F2\_0 & Accuracy & F1 \\
\midrule
1,000 & w10 & 0.555 & 0.910 & 0.806 & 0.922 & 0.955 \\
2,000 & w10 & 0.584 & 0.920 & 0.825 & 0.930 & 0.960 \\
5,000 & balanced & 0.606 & 0.935 & 0.843 & 0.936 & 0.963 \\
10,000 & balanced & 0.621 & 0.936 & 0.850 & 0.940 & 0.966 \\
\bottomrule
\end{tabular}
\end{table}

\subsection{So sánh biến thể}
Sau khi quét toàn bộ biến thể, cài đặt cổ điển được chọn là V6 với K bằng 10.000 và trọng lượng hạng w10. Bảng \ref{tab:nlp_variant_test} hiển thị so sánh thử nghiệm giữa các biến thể.

\begin{table}[H]
\centering
\caption{Các biến thể tình cảm NLP cổ điển trên dữ liệu thử nghiệm}
\label{tab:nlp_variant_test}
\begin{tabular}{lcccccc}
\toprule
Biến thể & K & Trọng số & Precision\_0 & Recall\_0 & F2\_0 & Coverage \\
\midrule
V0 & 10,000 & balanced & 0.638 & 0.937 & 0.856 & 0.870 \\
V1 & 10,000 & balanced & 0.643 & 0.932 & 0.855 & 0.871 \\
V2 & 5,000 & w10 & 0.647 & 0.926 & 0.852 & 0.874 \\
V3 & 10,000 & balanced & 0.626 & 0.923 & 0.843 & 0.868 \\
V4 & 10,000 & balanced & 0.635 & 0.921 & 0.845 & 0.872 \\
V5 & 10,000 & balanced & 0.653 & 0.933 & 0.859 & 0.878 \\
V6 & 10,000 & w10 & 0.677 & 0.932 & 0.867 & 0.885 \\
V7 & 5,000 & w10 & 0.657 & 0.927 & 0.857 & 0.880 \\
\bottomrule
\end{tabular}
\end{table}

V6 mang lại sự cân bằng thuận lợi cho mục tiêu của dự án. Nó duy trì khả năng thu hồi âm bản cao đồng thời cải thiện độ chính xác âm tính so với hầu hết các biến thể thay thế.

\subsection{Ngưỡng và hành vi không chắc chắn}
Đối với mẫu cổ điển V6 đã chọn, ngưỡng mặc định 0,40 và 0,60 mang lại phạm vi bao phủ 0,885 và thu hồi có chọn lọc\_0 0,956. Bảng \ref{tab:nlp_classic_threshold} hiển thị mức quét ngưỡng.

\begin{table}[H]
\centering
\caption{Quét ngưỡng cho mô hình tình cảm NLP cổ điển đã chọn}
\label{tab:nlp_classic_threshold}
\begin{tabular}{lcccc}
\toprule
Cặp ngưỡng & Coverage & Selective precision\_0 & Selective recall\_0 & Selective F2\_0 \\
\midrule
0.40 / 0.60 & 0.885 & 0.730 & 0.956 & 0.900 \\
0.45 / 0.60 & 0.891 & 0.708 & 0.956 & 0.894 \\
0.50 / 0.65 & 0.886 & 0.681 & 0.963 & 0.889 \\
0.55 / 0.70 & 0.881 & 0.646 & 0.970 & 0.882 \\
\bottomrule
\end{tabular}
\end{table}

\section{Kết quả cho phần mở rộng Transformer NLP}
\subsection{Số liệu máy biến áp}
Kết quả tinh chỉnh máy biến áp được lưu trong đầu ra mở rộng NLP. Bảng \ref{tab:nlp_transformer_metrics} tóm tắt các giá trị xác thực và kiểm tra.

\begin{table}[H]
\centering
\caption{Số liệu mở rộng Transformer NLP}
\label{tab:nlp_transformer_metrics}
\begin{tabular}{lcccccc}
\toprule
Tập dữ liệu & Accuracy & F1 & Precision\_0 & Recall\_0 & Coverage & Selective recall\_0 \\
\midrule
Validation & 0.968 & 0.982 & 0.821 & 0.854 & 0.908 & 0.862 \\
Test & 0.967 & 0.982 & 0.822 & 0.833 & 0.910 & 0.846 \\
\bottomrule
\end{tabular}
\end{table}

So với động cơ V6 cổ điển, máy biến áp cho độ chính xác tổng thể cao hơn và độ chính xác âm cao hơn ở ngưỡng cơ sở. Động cơ V6 cổ điển vẫn giữ mức thu hồi âm thô cao hơn nhiều.

\subsection{Quét ngưỡng máy biến áp}
Bảng \ref{tab:nlp_transformer_threshold} hiển thị hành vi chọn lọc của máy biến áp khi quét ngưỡng.

\begin{table}[H]
\centering
\caption{Quét ngưỡng máy biến áp trên dữ liệu thử nghiệm}
\label{tab:nlp_transformer_threshold}
\begin{tabular}{lcccc}
\toprule
Cặp ngưỡng & Coverage & Selective precision\_0 & Selective recall\_0 & Selective F2\_0 \\
\midrule
0.40 / 0.60 & 0.910 & 0.831 & 0.846 & 0.843 \\
0.45 / 0.60 & 0.911 & 0.827 & 0.847 & 0.843 \\
0.50 / 0.65 & 0.912 & 0.821 & 0.850 & 0.844 \\
0.55 / 0.70 & 0.911 & 0.818 & 0.855 & 0.847 \\
\bottomrule
\end{tabular}
\end{table}

Đường cong máy biến áp ổn định, có phạm vi bao phủ gần 0,91 và mức tăng vừa phải trong thu hồi có chọn lọc khi cặp ngưỡng di chuyển đến 0,55 và 0,70.

\subsection{So sánh trường hợp cứng}
Kết quả trường hợp cứng trong Bảng \ref{tab:nlp_hard_cases} so sánh các quyết định về đường cơ sở và máy biến áp cổ điển.

\begin{table}[H]
\centering
\caption{So sánh trường hợp khó giữa mô hình tình cảm cổ điển và biến đổi}
\label{tab:nlp_hard_cases}
\begin{tabular}{>{\raggedright\arraybackslash}p{4.1cm} >{\centering\arraybackslash}p{2.3cm} >{\centering\arraybackslash}p{2.3cm} >{\raggedright\arraybackslash}p{4.3cm}}
\toprule
Văn bản & Quyết định baseline & Quyết định transformer & Nhận xét chính \\
\midrule
not bad & Tiêu cực & Tích cực & Transformer diễn giải cụm thành ngữ theo hướng tích cực \\
not good & Tiêu cực & Tiêu cực & Hai mô hình cùng dự đoán theo hướng tiêu cực \\
good but late delivery & Tiêu cực & Tích cực & Transformer đặt trọng số cao hơn cho mệnh đề tích cực ở đầu câu \\
gr8 & Không chắc chắn & Không chắc chắn & Cả hai mô hình đều trả về không chắc chắn cho văn bản quá ngắn \\
thx & Không chắc chắn & Không chắc chắn & Cả hai mô hình đều trả về không chắc chắn cho văn bản quá ngắn \\
idk & Không chắc chắn & Không chắc chắn & Cả hai mô hình đều trả về không chắc chắn cho văn bản quá ngắn \\
\bottomrule
\end{tabular}
\end{table}

Sự so sánh này chỉ ra rằng mô hình Transformer giải quyết một số mẫu thành ngữ (ví dụ \textit{not bad}) hiệu quả hơn, trong khi đôi khi lại đánh giá quá cao các mệnh đề tích cực trong những lời phàn nàn có nhiều phân cực.

\section{Kết quả trích xuất vấn đề NLP}
\subsection{Kết quả đa nhãn tổng thể}
Bảng \ref{tab:nlp_issue_overall} so sánh hai thiết lập trích xuất vấn đề có thể tái tạo trong kho lưu trữ hiện tại.

\begin{table}[H]
\centering
\small
\setlength{\tabcolsep}{4pt}
\caption{Số liệu tổng thể về trích xuất vấn đề NLP}
\label{tab:nlp_issue_overall}
\begin{tabularx}{\textwidth}{
>{\raggedright\arraybackslash}X
>{\centering\arraybackslash}p{2.1cm}
>{\centering\arraybackslash}p{1.8cm}
>{\centering\arraybackslash}p{1.8cm}
>{\centering\arraybackslash}p{2.4cm}}
\toprule
Setup & Val micro-F1 & Test micro-F1 & Test macro-F1 & Test subset accuracy \\
\midrule
Word TF-IDF plus OVR Logistic & 0.759 & 0.765 & 0.381 & 0.606 \\
Word plus Char TF-IDF plus OVR Logistic & 0.764 & 0.763 & 0.408 & 0.603 \\
\bottomrule
\end{tabularx}
\end{table}

Việc biểu diễn từ cộng với ký tự mang lại macro-F1 cao hơn, trong khi thiết lập chỉ từ giữ độ chính xác của micro-F1 và tập hợp con cao hơn một chút trong lần chạy này.

\subsection{Cải thiện trên mỗi nhãn}
Bảng \ref{tab:nlp_issue_per_label_delta} hiển thị sự thay đổi F1 của thử nghiệm mỗi nhãn từ từ này sang từ khác cộng với thiết lập ký tự.

\begin{table}[H]
\centering
\caption{Thay đổi F1 kiểm tra trên mỗi nhãn trong trích xuất vấn đề}
\label{tab:nlp_issue_per_label_delta}
\begin{tabular}{lccc}
\toprule
Label & F1 word & F1 word plus char & Delta \\
\midrule
delivery\_shipping & 0.475 & 0.462 & -0.013 \\
redemption\_activation & 0.267 & 0.304 & +0.037 \\
product\_quality & 0.308 & 0.301 & -0.007 \\
customer\_service & 0.314 & 0.185 & -0.128 \\
refund\_return & 0.145 & 0.333 & +0.188 \\
usability & 0.462 & 0.476 & +0.014 \\
value\_price & 0.284 & 0.316 & +0.032 \\
fraud\_scam & 0.265 & 0.388 & +0.123 \\
other & 0.911 & 0.910 & -0.001 \\
\bottomrule
\end{tabular}
\end{table}

Các tính năng ký tự giúp một số nhãn thưa thớt như hoàn tiền\_return và lừa đảo\_scam, nhưng có thể giảm F1 trên một số nhãn. Điều này cho thấy sự đánh đổi về việc thu hồi độ chính xác theo nhãn.

\subsection{Chiến lược ngưỡng trong cài đặt nhiều nhãn}
Trích xuất vấn đề sử dụng ngưỡng điều chỉnh cho mỗi nhãn. Cấu hình ngưỡng khác nhau giữa thiết lập chỉ từ và thiết lập từ+char vì phân phối điểm thay đổi theo nhãn.

Hành vi này được mong đợi trong các tác vụ nhiều nhãn NLP trong đó mỗi nhãn có tần số trước và cường độ tín hiệu văn bản khác nhau.

\subsection{Tính tái lập của bước tạo template}
Bước tạo template gán nhãn issue hiện áp dụng cơ chế sắp xếp ID tự nhiên trước khi xuất CSV. Khóa sắp xếp được triển khai theo thứ tự: ID dạng số được sắp theo giá trị số nguyên, tiếp theo là ID không thuần số theo thứ tự từ vựng không phân biệt hoa thường, và cuối cùng là ID rỗng. Cách này loại bỏ lỗi thứ tự kiểu từ điển như \texttt{10} đứng trước \texttt{2}.

Hai test hồi quy được thêm tại \texttt{tests/test\_issue\_template\_sort.py}: một test đơn vị cho khóa sắp xếp và một test đầu-cuối xác nhận thứ tự output từ \texttt{cmd\_make\_template} trong \texttt{src/issue\_steps/steps.py}.

Cập nhật này cải thiện tính tái lập của dữ liệu chú thích và độ tin cậy khi bàn giao giữa các thành viên, đồng thời không làm thay đổi các chỉ số đánh giá mô hình đã báo cáo ở trên.


\subsection{C?p nh?t Level 3 cho classic pipeline}
Tr?ng th?i code hi?n t?i b? sung g?i Level 3 cho ph?n classic, g?m 4 ?i?m:
\begin{enumerate}
\item T?m class-weight theo t?ng label tr?n t?p validation (\texttt{none}, \texttt{balanced}, \texttt{w2}, \texttt{w5}, \texttt{w10}).
\item H? tr? calibration x?c su?t cho OVR Logistic Regression (\texttt{sigmoid} ho?c \texttt{isotonic}).
\item Xu?t b?o c?o ?? ?n ??nh ng??ng khi thay ??i delta \{-0.10, -0.05, 0.00, 0.05, 0.10\}.
\item Th?m l?nh ablation c? th? ch?y tr?c ti?p: \texttt{python -m src.nlp\_ext classic\_ablation}.
\end{enumerate}

\begin{table}[H]
\centering
\caption{Artefact Level 3 v? k?t qu? ch?nh t? l?n ch?y m?i nh?t}
\label{tab:nlp_level3_artifacts_vi}
\begin{tabular}{p{5.0cm} p{7.8cm}}
\toprule
Artefact & Ghi nh?n \\
\midrule
\texttt{results/issue\_steps/02\_class\_weight\_tuning.csv} & L?a ch?n theo label: 8 label d?ng \texttt{balanced}, label \texttt{other} d?ng \texttt{none}. \\
\texttt{results/issue\_steps/02\_threshold\_stability\_summary.md} & C? b?o c?o ?? ?n ??nh ng??ng cho c? calibrated v? uncalibrated tr?n val/test. \\
\texttt{results/nlp\_ext/syllabus\_upgrade/nlp\_ablation.csv} & ?? ch?y ablation V0, V2, V5, V6, V7 b?ng m?t command. \\
\texttt{results/nlp\_ext/syllabus\_upgrade/nlp\_ablation\_summary.md} & Trong l?n ch?y m?i nh?t, V7 cho trade-off negative-first t?t nh?t (recall\_0=0.869, precision\_0=0.636, F2\_0=0.810). \\
\bottomrule
\end{tabular}
\end{table}

\section{Kết quả cho Gói nâng cấp giáo trình}
\subsection{Đường cơ sở cổ điển cho các thuật toán khóa học cốt lõi}
Bảng \ref{tab:nlp_syllabus_bench} báo cáo số liệu kiểm tra cho khối cơ sở được thêm vào. Khối này sử dụng cùng một cách phân chia dữ liệu và không thêm nhãn mới.

\begin{table}[H]
\centering
\caption{Kết quả kiểm tra các cơ sở giáo trình NLP bổ sung}
\label{tab:nlp_syllabus_bench}
\begin{tabular}{lccc}
\toprule
Model & Precision\_0 & Recall\_0 & F2\_0 \\
\midrule
Logistic Regression (L2) & 0.703 & 0.913 & 0.862 \\
Multinomial Naive Bayes & 0.633 & 0.913 & 0.839 \\
Complement Naive Bayes & 0.510 & 0.952 & 0.811 \\
Perceptron & 0.684 & 0.828 & 0.795 \\
LinearSVM & 0.789 & 0.847 & 0.835 \\
SGD log-loss & 0.616 & 0.921 & 0.838 \\
SVD semantic Logistic Regression & 0.501 & 0.920 & 0.788 \\
FFNN MLP on SVD vectors & 0.799 & 0.730 & 0.742 \\
\bottomrule
\end{tabular}
\end{table}

Những kết quả này cho thấy sự khác biệt rõ ràng về hành vi. Complement Naive Bayes có mức thu hồi âm cao nhất là 0,952, nhưng độ chính xác thấp ở mức 0,510. Hồi quy logistic cho F2\_0 tốt nhất ở mức 0,862 với độ chính xác cao hơn. LinearSVM có độ chính xác cao nhất\_0 trong số các mô hình thu hồi mạnh mẽ, trong khi tính năng ghi nhật ký SGD mang lại khả năng thu hồi cao với độ chính xác thấp hơn. Điều này hỗ trợ một cuộc thảo luận đánh đổi mạnh mẽ giữa các thuật toán khóa học cốt lõi.

\subsection{Hành vi chọn lọc trong các đường cơ sở bổ sung}
Đối với các mô hình xác suất trong khối được thêm vào, áp dụng cùng một dải không chắc chắn 0,40 và 0,60. Hồi quy logistic cho phạm vi kiểm tra 0,887, độ chính xác chọn lọc\_0 0,763 và thu hồi có chọn lọc\_0 0,946. Phần bổ sung Naive Bayes cho phạm vi kiểm tra 0,900 và thu hồi có chọn lọc\_0 0,963, nhưng độ chính xác chọn lọc\_0 chỉ là 0,545. Điều này xác nhận rằng khối được thêm vào phù hợp với chính sách của dự án và duy trì sự kiểm soát cân bằng minh bạch.

\subsection{Kết quả mô hình ngôn ngữ N-Gram}
Bảng \ref{tab:nlp_ngram_results} báo cáo mức độ phức tạp của các mô hình unigram và bigram với add-k bằng 1,0.

\begin{table}[H]
\centering
\caption{Sự bối rối của khối mô hình ngôn ngữ N-gram}
\label{tab:nlp_ngram_results}
\begin{tabular}{lcc}
\toprule
Model & Validation perplexity & Test perplexity \\
\midrule
Unigram add-k & 376.51 & 378.02 \\
Bigram add-k & 358.59 & 363.08 \\
\bottomrule
\end{tabular}
\end{table}

Bigram mang lại độ phức tạp thấp hơn unigram trên cả hai phần tách. Điều này được mong đợi vì bigram sử dụng ngữ cảnh cục bộ từ mã thông báo trước đó. Khoảng cách không lớn vì văn bản ôn tập ngắn và ồn ào nhưng ổn định và có ý nghĩa đối với chủ đề mô hình ngôn ngữ.

\subsection{Các tạo tác RNN, MLM và LLM}
Bảng \ref{tab:nlp_extra_artifacts} tóm tắt các tạo phẩm giáo trình bổ sung hiện đang được triển khai và có thể tái tạo ở trạng thái kho lưu trữ này.

\begin{table}[H]
\centering
\caption{Các tạo phẩm giáo trình NLP bổ sung}
\label{tab:nlp_extra_artifacts}
\begin{tabular}{lccc}
\toprule
Tạo tác & Chỉ số & Giá trị & Ghi chú \\
\midrule
RNN/LSTM baseline & recall\_0 (test) & 0.911 & recall cao, precision còn hạn chế \\
RNN/LSTM baseline & F2\_0 (test) & 0.615 & hỗ trợ bao phủ chủ đề mô hình chuỗi \\
MLM probe & hit@10 & 0.600 & 3/5 mẫu probe khớp với tập token kỳ vọng \\
LLM prompt baseline & recall\_0 (test) & 0.513 & đường cơ sở proxy theo prompt với đầu ra xác định \\
LLM prompt baseline & F2\_0 (test) & 0.499 & thấp hơn các baseline phân loại, giữ lại để thể hiện độ bao phủ học phần \\
\bottomrule
\end{tabular}
\end{table}

\subsection{Ma trận phù hợp với khóa học được cập nhật}
Bảng \ref{tab:nlp_course_fit} hiển thị ma trận phạm vi chủ đề được cập nhật sau khối nâng cấp mới.

\begin{table}[H]
\centering
\caption{Phạm vi chủ đề được cập nhật sau khi nâng cấp không có dữ liệu mới}
\label{tab:nlp_course_fit}
\begin{tabular}{lcc}
\toprule
Nhóm chủ đề & Tỷ lệ bao phủ & Mức độ \\
\midrule
Xác suất & 100 & Cao \\
Biểu thức chính quy & 50 & Một phần \\
Các thước đo đánh giá & 100 & Cao \\
Mô hình ngôn ngữ N-gram & 100 & Cao \\
Naive Bayes & 100 & Cao \\
Perceptron & 100 & Cao \\
Hồi quy logistic & 100 & Cao \\
Mạng nơ-ron truyền thẳng & 100 & Cao \\
PyTorch cơ bản & 100 & Cao \\
Ngữ nghĩa vector và embedding & 80 & Trung bình-cao \\
Mạng nơ-ron hồi tiếp & 100 & Cao \\
Transformer và LLM & 100 & Cao \\
Mô hình ngôn ngữ che khuyết & 100 & Cao \\
Ứng dụng của LLM & 90 & Cao \\
\bottomrule
\end{tabular}
\end{table}

Điểm bao quát tổng thể ước tính là 94,3\%, được tính trực tiếp từ các tạo phẩm được tạo trong thư mục đầu ra nâng cấp giáo trình.

\section{Thảo luận về máy biến áp cổ điển và máy biến áp}
NLP thưa thớt cổ điển và NLP dựa trên máy biến áp thể hiện các cấu hình hiệu suất khác biệt. Trong cài đặt này, động cơ V6 cổ điển duy trì khả năng thu hồi âm cao và thu hồi có chọn lọc cao, trong khi mẫu máy biến áp cung cấp độ chính xác tổng thể cao hơn và độ chính xác âm cao hơn.

Đối với các kịch bản phân loại ưu tiên phát hiện các đánh giá bất lợi, V6 cổ điển phù hợp chặt chẽ hơn với chính sách quyết định. Đối với sắc thái ngữ nghĩa và cách diễn đạt thành ngữ, đầu ra máy biến áp cung cấp giá trị bổ sung khi kết hợp với điều khiển ngưỡng.

Do đó, chiến lược kết hợp là hợp lý: sử dụng mô hình cổ điển làm lớp phân loại chính và áp dụng tính điểm biến áp làm bước phân tích phụ cho các trường hợp không rõ ràng.

\section{Phân tích lỗi và các chế độ sai hỏng}
Bảng \ref{tab:nlp_error_cases} trình bày các trường hợp lỗi hoặc rủi ro cụ thể kèm theo giải thích cụ thể cho từng mô hình. Điều này mở rộng việc đánh giá ra ngoài các số liệu tổng hợp.

\begin{table}[H]
\centering
\caption{Các trường hợp lỗi và rủi ro cụ thể kèm theo lời giải thích}
\label{tab:nlp_error_cases}
\begin{tabular}{>{\raggedright\arraybackslash}p{2.9cm} >{\centering\arraybackslash}p{2.0cm} >{\raggedright\arraybackslash}p{2.7cm} >{\raggedright\arraybackslash}p{4.4cm} >{\raggedright\arraybackslash}p{2.6cm}}
\toprule
Văn bản & Kỳ vọng & Đầu ra mô hình & Nguyên nhân sai/rủi ro & Biện pháp khắc phục trong dự án \\
\midrule
not bad & Tích cực & classic: Tiêu cực & Biểu diễn thưa đặt trọng số quá cao cho token bad và chưa mô hình hóa tốt phủ định thành ngữ & Duy trì so sánh với transformer và cơ chế phân xử theo dải không chắc chắn khi độ tin cậy thấp \\
good but late delivery & Tiêu cực hoặc cần gắn cờ vấn đề & transformer: Tích cực & Chấm điểm theo ngữ cảnh đặt trọng số cao ở mệnh đề tích cực đầu câu và thiếu nhấn mạnh phần phàn nàn cuối câu & Giữ mô hình classic theo chính sách negative-first để định tuyến và dùng mô-đun issue extraction để gắn nhãn khiếu nại \\
gr8 / thx / idk & Không chắc chắn & both: Không chắc chắn & Văn bản quá ngắn mang tín hiệu cảm xúc yếu và độ mơ hồ cao & Duy trì luật fallback too\_short để tránh gán nhãn cứng cưỡng bức \\
refund\_return (nhãn hiếm) & Cần phát hiện đúng vấn đề hiếm & recall của mô hình issue dao động & Số mẫu hỗ trợ thấp làm biên quyết định kém ổn định và nhạy với ngưỡng & Áp dụng biến thể char n-gram và tinh chỉnh ngưỡng theo nhãn; ưu tiên bổ sung gán nhãn cho nhóm hiếm \\
\bottomrule
\end{tabular}
\end{table}

Các kiểu lỗi phổ biến là phân cực hỗn hợp với các điểm đánh dấu tương phản, văn bản có ít thông tin rất ngắn và tính không ổn định của nhãn thưa trong quá trình trích xuất nhiều nhãn. Các mẫu này được theo dõi và giải quyết rõ ràng thông qua các biện pháp kiểm soát phiếu trắng và theo từng nhãn.

\section{Giao thức trình diễn}
Giao thức trình diễn trực tiếp được thiết kế cho cửa sổ trình bày kéo dài 10 phút với hoạt động ổn định và có thể kiểm tra được.

\subsection{Luồng trình diễn (10 phút bao gồm cả phần hỏi đáp)}
\begin{enumerate}
\item Kiểm tra độ tái lập: chạy \texttt{python -m pytest -q} và hiển thị đã vượt qua các bài kiểm tra khói.
\item Chạy bản demo CLI cảm tính với đầu vào tích cực, tiêu cực và kiểu không chắc chắn.
\item Chạy dự đoán trích xuất vấn đề cho văn bản giống như khiếu nại.
\item Chạy bản demo máy biến áp cho một cụm từ ngắn mơ hồ.
\item Hiển thị bảng điểm thống nhất và giải thích sự cân bằng giữa các mô hình.
\end{enumerate}

\subsection{Các lệnh được đề xuất}
\begin{verbatim}
python -m pytest -q
python demo.py "great product and fast delivery" "terrible quality, broke after two days" "idk if this is worth the price"
python -m src.issue_steps predict --text "good but slow delivery"
python demo_transformer.py "not bad at all"
python scripts/build_scoreboard.py
\end{verbatim}

\subsection{Hành vi demo dự kiến}
Hành vi dự kiến được chỉ định trong \texttt{docs/expected\_outputs.md}. Các đánh giá tích cực và tiêu cực rõ ràng phải chính xác theo hướng, trong khi thông tin thấp hoặc tín hiệu hỗn hợp có thể trả về \texttt{Không chắc chắn}. Giao thức này hỗ trợ khả năng diễn giải và ổn định hoạt động.

\section{Ma trận truy vết tạo tác}
Bảng \ref{tab:nlp_project_evidence} ánh xạ từng khía cạnh dự án tới bằng chứng cụ thể trong kho lưu trữ này.

\begin{table}[H]
\centering
\caption{Kích thước dự án và truy xuất nguồn gốc tạo tác}
\label{tab:nlp_project_evidence}
\begin{tabular}{>{\raggedright\arraybackslash}p{2.5cm} >{\raggedright\arraybackslash}p{5.2cm} >{\raggedright\arraybackslash}p{4.7cm}}
\toprule
Tiêu chí & Nội dung đã triển khai & Bằng chứng trong kho mã \\
\midrule
Nội dung trình bày & Mô tả pipeline rõ ràng, lệnh chạy có thể tái lập, luồng demo có thể kiểm chứng & Các lệnh trong README, \texttt{src/run\_all.py}, \texttt{demo.py}, \texttt{demo\_transformer.py}, \texttt{demo\_app.py} \\
Ý nghĩa dự án & Bài toán thực tế theo ngữ cảnh triage hỗ trợ, chính sách negative-first và nhãn hành động mức issue & \texttt{src/dm2\_steps/common.py}, các tạo tác uncertainty, CSV so sánh ca khó \\
Quy trình làm việc & Phân công module rõ ràng, quy tắc review chéo, cổng kiểm soát chất lượng và kiểm tra thứ tự template issue theo tính xác định & \texttt{docs/contribution\_matrix.md}, \texttt{.github/pull\_request\_template.md}, smoke tests, \texttt{tests/test\_issue\_template\_sort.py} \\
Chất lượng kết quả & So sánh nhiều mô hình, chỉ số selective, phân tích theo nhãn và giải thích ca lỗi & \texttt{results/scoreboard/*}, \texttt{results/nlp\_ext/syllabus\_upgrade/*}, phần phân tích lỗi \\
Chất lượng demo & Có cả đường chạy CLI và UI, đầu vào kiểm thử xác định, tiêu chí hành vi kỳ vọng rõ ràng & \texttt{docs/demo\_inputs.txt}, \texttt{docs/expected\_outputs.md}, demo Streamlit và CLI \\
\bottomrule
\end{tabular}
\end{table}

\section{Mức độ phù hợp học thuật với học phần NLP}
Dự án này hiện hiển thị nhiều ý tưởng NLP trong một hệ thống. Nó bao gồm chuẩn hóa văn bản, ánh xạ từ vựng, mã hóa thưa thớt n-gram, xử lý phạm vi phủ định, tín hiệu cấp mệnh đề, phân loại văn bản nhiều nhãn, hiệu chỉnh ngưỡng và tinh chỉnh biến áp. Nó cũng bao gồm Naive Bayes, Perceptron, đường cơ sở ngữ nghĩa chuyển tiếp nguồn cấp dữ liệu rõ ràng, tạo phẩm mô hình ngôn ngữ N-gram, đường cơ sở RNN/LSTM, thăm dò MLM và tạo phẩm ứng dụng kiểu lời nhắc LLM.

Dự án vẫn có thể sử dụng trực tiếp cho thử nghiệm dựa trên khóa học. Nó cung cấp các tạo phẩm đã lưu, số liệu có thể tái tạo, bản ghi quy trình làm việc chú thích và tập lệnh suy luận có thể triển khai. Nó cũng bao gồm một ma trận bao phủ chủ đề phù hợp cho việc bảo vệ miệng.

Đối với việc học trong khóa học, phiên bản cập nhật cung cấp bằng chứng mạnh mẽ hơn về sự cân bằng giữa các phương pháp tiếp cận NLP cổ điển và sâu sắc theo một mục tiêu kinh doanh và một tập dữ liệu cố định.

\section{Hạn chế và công việc trong tương lai}
Hạn chế chính vẫn là sự mất cân bằng dữ liệu, đặc biệt là về tâm lý tiêu cực và các nhãn vấn đề hiếm gặp. Điều này ảnh hưởng đến độ chính xác và độ ổn định trong một số mô hình.

Một hạn chế khác là độ sâu không đồng đều giữa các chủ đề. Mặc dù mức độ bao phủ cao nhưng một số chủ đề được thể hiện bằng các tạo phẩm nhẹ (ví dụ: các biểu thức chính quy ở mức độ bao phủ 50\% và các đầu dò MLM/LLM nhỏ gọn).

Công việc trong tương lai có thể tập trung vào chiều sâu chất lượng chứ không chỉ sự hiện diện của chủ đề. Các bước tiếp theo là: tạo tác cắt bỏ hoàn toàn cho các khối phủ định/char/từ vựng, phân nhóm lỗi phong phú hơn cho các phủ định sai, phân tích tiền xử lý tập trung vào biểu thức chính quy mạnh mẽ hơn và công cụ demo chặt chẽ hơn liên kết từng dự đoán với mô-đun đóng góp của nó.

\section{Phần kết luận}
Công việc này trình bày một quy trình đánh giá-hiểu biết hoàn chỉnh bao gồm phân loại tình cảm, trích xuất vấn đề và mở rộng dựa trên máy biến áp. Mô hình thưa thớt nhận biết ngữ cảnh cổ điển cung cấp hành vi tiêu cực đầu tiên mạnh mẽ, trong khi phần mở rộng biến áp cải thiện độ chính xác và độ chắc chắn theo ngữ cảnh trong các trường hợp được chọn. Trích xuất vấn đề đa nhãn chuyển đổi kết quả cảm tính thành các danh mục vấn đề có liên quan đến hoạt động. Khối mở rộng giáo trình bao gồm các đường cơ sở cổ điển, mô hình ngôn ngữ N-gram, RNN/LSTM, thăm dò MLM và các tạo phẩm nhắc nhở LLM mà không đưa ra dữ liệu được gắn nhãn mới. Với bộ tạo tác hiện tại, dự án thể hiện khả năng tái tạo và độ bao phủ giáo trình ước tính là 94,3\%.

\clearpage
\section*{Tài liệu tham khảo}
\addcontentsline{toc}{section}{Tài liệu tham khảo}
\begin{thebibliography}{99}

\bibitem{jurafsky2026}
Daniel Jurafsky and James H. Martin.
\textit{Speech and Language Processing} (3rd ed. draft).
Online manuscript, January 2026.
\newline\texttt{https://web.stanford.edu/~jurafsky/slp3/}

\bibitem{manning2008ir}
Christopher D. Manning, Prabhakar Raghavan, and Hinrich Schutze.
\textit{Introduction to Information Retrieval}.
Cambridge University Press, 2008.
\newline\texttt{https://nlp.stanford.edu/IR-book/}

\bibitem{turney2002}
Peter D. Turney.
Thumbs up or thumbs down? Semantic orientation applied to unsupervised classification of reviews.
In \textit{Proceedings of ACL}, 2002.
\newline\texttt{https://aclanthology.org/P02-1053/}

\bibitem{wang2012}
Sida I. Wang and Christopher D. Manning.
Baselines and Bigrams: Simple, Good Sentiment and Topic Classification.
In \textit{Proceedings of ACL}, 2012.
\newline\texttt{https://aclanthology.org/P12-2018/}

\bibitem{wiegand2010}
Michael Wiegand, Alexandra Balahur, Benjamin Roth, Dietrich Klakow, and Andres Montoyo.
A Survey on the Role of Negation in Sentiment Analysis.
In \textit{Proceedings of the Workshop on Negation and Speculation in Natural Language Processing}, 2010.
\newline\texttt{https://aclanthology.org/W10-3111/}

\bibitem{joulin2017}
Armand Joulin, Edouard Grave, Piotr Bojanowski, and Tomas Mikolov.
Bag of Tricks for Efficient Text Classification.
In \textit{Proceedings of EACL}, 2017.
\newline\texttt{https://aclanthology.org/E17-2068/}

\bibitem{nielsen2011}
Finn Arup Nielsen.
A new ANEW: Evaluation of a word list for sentiment analysis in microblogs.
\textit{arXiv preprint arXiv:1103.2903}, 2011.
\newline\texttt{https://arxiv.org/abs/1103.2903}

\bibitem{forman2003}
George Forman.
An extensive empirical study of feature selection metrics for text classification.
\textit{Journal of Machine Learning Research}, 3:1289--1305, 2003.
\newline\texttt{https://www.jmlr.org/papers/v3/forman03a.html}

\bibitem{rifkin2004}
Ryan Rifkin and Aldebaro Klautau.
In defense of one-vs-all classification.
\textit{Journal of Machine Learning Research}, 5:101--141, 2004.
\newline\texttt{https://www.jmlr.org/papers/v5/rifkin04a.html}

\bibitem{tsoumakas2011}
Grigorios Tsoumakas, Ioannis Katakis, and Ioannis Vlahavas.
Random k-Labelsets for multilabel classification.
\textit{IEEE Transactions on Knowledge and Data Engineering}, 23(7):1079--1089, 2011.
\newline\texttt{https://doi.org/10.1109/TKDE.2010.164}

\bibitem{zhang2014}
Min-Ling Zhang and Zhi-Hua Zhou.
A review on multi-label learning algorithms.
\textit{IEEE Transactions on Knowledge and Data Engineering}, 26(8):1819--1837, 2014.
\newline\texttt{https://doi.org/10.1109/TKDE.2013.39}

\bibitem{vaswani2017}
Ashish Vaswani, Noam Shazeer, Niki Parmar, Jakob Uszkoreit, Llion Jones, Aidan N. Gomez, Lukasz Kaiser, and Illia Polosukhin.
Attention Is All You Need.
In \textit{Advances in Neural Information Processing Systems 30}, 2017.
\newline\texttt{https://arxiv.org/abs/1706.03762}

\bibitem{devlin2019}
Jacob Devlin, Ming-Wei Chang, Kenton Lee, and Kristina Toutanova.
BERT: Pre-training of Deep Bidirectional Transformers for Language Understanding.
In \textit{Proceedings of NAACL-HLT}, 2019.
\newline\texttt{https://aclanthology.org/N19-1423/}

\bibitem{saito2015}
Takaya Saito and Marc Rehmsmeier.
The precision-recall plot is more informative than the ROC plot when evaluating binary classifiers on imbalanced datasets.
\textit{PLOS ONE}, 10(3):e0118432, 2015.
\newline\texttt{https://doi.org/10.1371/journal.pone.0118432}

\bibitem{geifman2017}
Yonatan Geifman and Ran El-Yaniv.
Selective Classification for Deep Neural Networks.
\textit{arXiv preprint arXiv:1705.08500}, 2017.
\newline\texttt{https://arxiv.org/abs/1705.08500}

\end{thebibliography}

\clearpage
\section*{Phụ lục A: Số liệu tình cảm}
\addcontentsline{toc}{section}{Phụ lục A: Hình kết quả cảm xúc}

\begin{figure}[H]
\centering
\includegraphics[width=0.82\textwidth]{\detokenize{results/dm2_steps/01_rating_distribution.png}}
\caption{Phân phối xếp hạng trong dữ liệu nguồn}
\end{figure}

\begin{figure}[H]
\centering
\includegraphics[width=0.82\textwidth]{\detokenize{results/dm2_steps/06_f1_vs_k.png}}
\caption{Xác thực F1 bằng Chi-square top K}
\end{figure}

\begin{figure}[H]
\centering
\includegraphics[width=0.82\textwidth]{\detokenize{results/dm2_steps/06_recall0_vs_k.png}}
\caption{Thu hồi xác thực cho lớp âm theo top K}
\end{figure}

\begin{figure}[H]
\centering
\includegraphics[width=0.82\textwidth]{\detokenize{results/dm2_steps/06b_recall0_by_variant.png}}
\caption{Kiểm tra việc thu hồi lớp phủ định bằng biến thể NLP cổ điển}
\end{figure}

\begin{figure}[H]
\centering
\includegraphics[width=0.82\textwidth]{\detokenize{results/dm2_steps/06b_precision0_by_variant.png}}
\caption{Kiểm tra độ chính xác cho lớp phủ định bằng biến thể NLP cổ điển}
\end{figure}

\begin{figure}[H]
\centering
\includegraphics[width=0.82\textwidth]{\detokenize{results/dm2_steps/09_prob_hist_test_strong.png}}
\caption{Biểu đồ xác suất mô hình cổ điển trên tập kiểm tra mạnh}
\end{figure}

\begin{figure}[H]
\centering
\includegraphics[width=0.82\textwidth]{\detokenize{results/dm2_steps/10_threshold_tradeoff.png}}
\caption{Sự đánh đổi ngưỡng mô hình cổ điển cho mức độ phù hợp và thu hồi}
\end{figure}

\begin{figure}[H]
\centering
\includegraphics[width=0.82\textwidth]{\detokenize{results/nlp_ext/nlp_confusion_matrix.png}}
\caption{Ma trận nhầm lẫn máy biến áp trên bộ thử nghiệm}
\end{figure}

\begin{figure}[H]
\centering
\includegraphics[width=0.82\textwidth]{\detokenize{results/nlp_ext/nlp_prob_hist_test.png}}
\caption{Biểu đồ xác suất máy biến áp trên tập thử nghiệm}
\end{figure}

\begin{figure}[H]
\centering
\includegraphics[width=0.82\textwidth]{\detokenize{results/nlp_ext/nlp_threshold_tradeoff.png}}
\caption{Sự đánh đổi ngưỡng máy biến áp để phủ sóng và thu hồi}
\end{figure}

\clearpage
\section*{Phụ lục B: Số liệu trích xuất vấn đề}
\addcontentsline{toc}{section}{Phụ lục B: Hình kết quả trích xuất vấn đề}

\begin{figure}[H]
\centering
\includegraphics[width=0.82\textwidth]{\detokenize{results/issue_steps/02_label_frequency.png}}
\caption{Tần suất phát hành nhãn trong bộ chú thích thủ công}
\end{figure}

\begin{figure}[H]
\centering
\includegraphics[width=0.82\textwidth]{\detokenize{results/issue_steps/02_per_label_f1.png}}
\caption{Mỗi nhãn F1 cho mô hình trích xuất vấn đề từ}
\end{figure}

\begin{figure}[H]
\centering
\includegraphics[width=0.82\textwidth]{\detokenize{results/issue_steps_char_demo/02_per_label_f1.png}}
\caption{Mỗi nhãn F1 cho mô hình trích xuất vấn đề từ cộng với char}
\end{figure}

\clearpage
\section*{Phụ lục C: Số liệu nâng cấp giáo trình}
\addcontentsline{toc}{section}{Phụ lục C: Hình kết quả nâng cấp học phần}

\begin{figure}[H]
\centering
\includegraphics[width=0.82\textwidth]{\detokenize{results/nlp_ext/syllabus_upgrade/nlp_syllabus_bench_recall0.png}}
\caption{Nhớ lại bài kiểm tra\_0 trên các cơ sở giáo trình được bổ sung}
\end{figure}

\begin{figure}[H]
\centering
\includegraphics[width=0.82\textwidth]{\detokenize{results/nlp_ext/syllabus_upgrade/nlp_syllabus_bench_f2.png}}
\caption{Kiểm tra F2\_0 trên các cơ sở giáo trình bổ sung}
\end{figure}


\begin{figure}[H]
\centering
\includegraphics[width=0.82\textwidth]{\detokenize{results/nlp_ext/syllabus_upgrade/nlp_ablation_f2.png}}
\caption{K?t qu? ablation classic: F2\_0 tr?n t?p test theo t?ng bi?n th?}
\end{figure}

\end{document}
